\section{Verification Ladder Specification}
\label{app:ladder}

The ladder is the finite-state machine of \S\ref{sec:ladder}.
This appendix supplies the per-rung specification used by the
Source Analyzer (\S\ref{sec:pipeline}) and the Aligner.

\begin{table}[!ht]
\centering
\small
\begin{tabular}{lp{2.4cm}p{4.0cm}p{4.5cm}}
\toprule
\textbf{Rung} & \textbf{Required artefact} & \textbf{Promotion condition} & \textbf{PRISMA / GRADE floor} \\
\midrule
\textsc{unverified-external}
  & --
  & Claim recognised as needing a source.
  & --- \\
\textsc{needs-research}
  & Search query in \texttt{research-protocol.md}
  & A candidate URL or DOI has been identified.
  & --- \\
\textsc{lit-retrieved}
  & URL/DOI that resolves
  & A Source Analyzer reads the source and confirms claim support, with quoted snippet.
  & PRISMA: Identification \\
\textsc{ai-confirmed}
  & Quoted snippet in \texttt{sources.md} plus \texttt{fair2r:Source} entity
  & A human reads the source and confirms the inference from source to claim.
  & PRISMA: Screening; GRADE floor: \textsc{low} \\
\textsc{lit-read}
  & Vendored source in \texttt{doc/sources/} (when licensable) or human signature in logbook
  & ---
  & PRISMA: Included; GRADE floor: \textsc{moderate} \\
\bottomrule
\end{tabular}
\caption{Verification ladder, with promotion condition and crosswalk
floor. The crosswalk is the floor each rung permits, not a guarantee of
the GRADE ceiling.}
\label{tab:ladder}
\end{table}

\paragraph{Retraction.}
A claim is retracted by adding a \texttt{prov:Invalidation} activity
that names the original \texttt{fair2r:Claim} as
\texttt{prov:invalidated}; the original triple is not deleted. A
retracted claim cannot be invoked by the manuscript and is rendered
visibly retracted on the public site.

\paragraph{Per-rung budget.}
At submission time, the manuscript reports the fraction of claims at
each rung. A typical first submission carries a single-digit
percentage at \textsc{lit-retrieved}, a moderate fraction at
\textsc{ai-confirmed}, and the load-bearing remainder at
\textsc{lit-read}. The ratio is itself a disclosure under the sixth
practice (\S\ref{sec:eight}, item~6).
